\section{Uniform phase stability}
\begin{proposition}[Uniform phase law and exact pairing]\label{v6:prop:allphase}\label{prop:T5b:symmetries}
For $\phi\in\mathbb R$, let $\mathcal K_{\tau,\phi}$ have kernel
\[
 \frac{\sin(2\pi\tau x+\phi)-\sin(2\pi\tau y+\phi)}{\pi(x-y)}.
\]
For every real Lipschitz $f$ with $f(0)=0$ and every $\tau>0$,
\[
 |\tr f(\mathcal K_{\tau,\phi})-\tr f(\mathcal K_\tau)|
 \le4b(1/2)\operatorname{Lip}(f),
\]
where $b$ is defined in \eqref{rev5:eq:strip-function}.
Consequently Theorem~\ref{T10:conj:law} holds uniformly in $\phi$,
with the same $\varkappa(f)$ and remainder constant
$C_{\mathrm{tr}}+4b(1/2)$.  Moreover
\begin{equation}\label{v6:eq:phasesym}
 \begin{gathered}
 \mathcal R\mathcal K_{\tau,\phi}\mathcal R=\mathcal K_{\tau,-\phi},\qquad
 \mathcal K_{\tau,\phi+\pi}=-\mathcal K_{\tau,\phi},\qquad
 \mathcal K_{\tau,\phi+2\pi}=\mathcal K_{\tau,\phi},\\
 \tr\mathcal K_{\tau,\phi}=\frac2\pi\sin(2\pi\tau)\cos\phi.
 \end{gathered}
\end{equation}
At $\phi=\pi/2\pmod\pi$ its spectrum has exact $\pm$ pairing with
multiplicity, and every odd Lipschitz test has zero trace.  Each eigenvector
for a nonzero eigenvalue has even and odd parts of equal norm.
\end{proposition}
\begin{proof}
The operator is unitarily equivalent to the compression of $\mathcal G$
to $(-\tau,\tau)+\phi/(2\pi)$.  Reduce $\phi$ modulo $2\pi$ so that
$|\phi/(2\pi)|\le1/2$, and apply Lemmas~\ref{rev5:lem:strips}
and~\ref{rev5:lem:trace-lip}.  The trace law follows since
$\operatorname{Lip}(f)\le\|f'\|_\infty$.
Reflection and phase shifts in the kernel give the first three identities;
integration of its diagonal gives the fourth.  At the indicated phases,
reflection anticommutes with the operator and therefore pairs eigenvalues
of opposite signs.  It exchanges the signs, rather than defining commuting
parity sectors there.  If $\mathcal K_{\tau,\phi}u=\lambda u$ with
$\lambda\ne0$, then $u$ and $\mathcal Ru$ belong to distinct eigenspaces
and are orthogonal.  Thus the two reflection components of $u$ have
squared norm $\|u\|^2/2$.
\end{proof}
At the centered phase $\phi=0$, neither parity spectrum is sign-symmetric,
for any $\tau>0$.  Indeed its trace in \eqref{eq:2.8} is nonzero because
$\Si(t)>|\sin t|$ for $t>0$.  On $(0,\pi]$, differentiate
$\Si(t)-\sin t$ and use $\sin t-t\cos t>0$.
On $[\pi,2\pi]$, use monotonicity and
$\Si(2\pi)\ge\pi/2-1/\pi>1$; for $t\ge2\pi$ the same inequality
follows from $|\Si(t)-\pi/2|\le2/t$.
A sign-symmetric trace-class spectrum would have zero trace.
The full centered spectrum is also asymmetric whenever $\sin(2\pi\tau)\ne0$.


